Our evaluation has four parts: a corpus of long documents split into length buckets, three summarizers, a human consistency judgement for every summary, and two automatic metrics. We describe each and then the two comparisons we draw between the human and the automatic view.

\subsection{Documents and Length Buckets}
The corpus consists of 200 long documents in English: feature articles from news magazines and reports published by public institutions. We chose these two sources because they are long, written for a general reader, and dense with the named entities, quantities and dates on which summarizers most often go wrong. We measured each document in subword tokens of the summarizers' tokenizer and assigned it to one of three buckets: short documents under 3,000 tokens, medium documents between 3,000 and 6,000 tokens, and long documents above 6,000 tokens. The short and medium buckets hold 70 documents each and the long bucket holds 60. The buckets are fixed before any summary is scored, and every analysis below is reported for the whole corpus and for each bucket separately. We do not claim that 6,000 tokens is a natural boundary; it is where the longest third of our corpus begins.

\subsection{Summarizers}
We compare three systems that span the extractive to abstractive range. Lead-3 takes the first three sentences of the document. It cannot introduce a fact that is absent from the source, but its sentences can lose the context that made them accurate, and its quality depends on how much a document front-loads its content \citep{grenander2019leadbias}. Bart-large is the BART sequence-to-sequence model fine-tuned for summarization \citep{lewis2020bart}; for inputs longer than its context we summarize the leading window, which is the common practice for this model. Bart-large-rl starts from the same checkpoint and is further trained with a reward from an entailment model that scores the summary against the source, in the spirit of reinforcement learning from textual entailment feedback \citep{roit2023rlef}. Contrastive and entity-level objectives pursue the same goal by other means \citep{cao2021cliff,nan2021entity}; we include one reward-trained model because it is the setting in which a metric could most easily be gamed: a model optimized towards an entailment signal may learn what the metric rewards rather than what annotators accept. Each system produced one summary per document, giving three summaries for each of the 200 documents.

\subsection{Human Judgement}
Annotators judged each summary against its full source document and answered one question: is every statement in the summary supported by the document? A summary with any unsupported statement is inconsistent. We use this binary judgement because it is the one a practitioner needs when deciding whether a summary can be shown to a reader, and because finer-grained schemes, although more informative, are costly on long sources \citep{krishna2023longeval,liu2023rose}. Two annotators with a background in journalism judged every summary independently. They saw the summaries in random order without the system names, and the three summaries of a document never appeared next to each other. Disagreements were resolved by a third annotator who saw both judgements and the source. Before the main round, the annotators judged a small calibration set of summaries of documents outside the corpus and discussed every disagreement, which settled how to treat paraphrase, rounding of quantities and statements that combine facts from distant parts of the document. Those rules were written down and applied unchanged in the main round. We report the agreement of the two first-round annotators as Cohen's kappa, and for every system and bucket the share of summaries judged consistent after adjudication, which we call the human consistency rate.

The judgement is deliberately coarse. It does not distinguish an unsupported adjective from an invented event, and it counts a summary with several errors the same as one with a single error. This matches how the automatic metrics are used in practice, as a single score per summary, and it keeps the annotation of long documents affordable.

\subsection{Automatic Metrics}
SummaC splits the source and the summary into sentences, scores every source-summary sentence pair with a natural language inference model trained on MNLI and related data \citep{laban2022summac,williams2018mnli}, and aggregates the pair scores into one summary score. We use the convolutional variant with its released default settings. QAFactEval extracts answer spans from the summary, generates a question for each, answers the questions against the source, and compares the answers \citep{fabbri2022qafacteval}. We use its released configuration. Both metrics return one score per summary, higher meaning more consistent. Neither was modified for long inputs: SummaC compares every summary sentence with every source sentence regardless of length, while QAFactEval retrieves the passage that answers each question. Their different ways of locating evidence in a long document are the reason we expect them to react differently to length. Both metrics read the full source text of every document, so any change with length reflects how they use the evidence rather than what they are shown. We scored all summaries of a document in the same run, and the metric scores were computed before the human labels were joined to them.

\subsection{Analysis}
We compare the metrics with the human judgement at two levels. At the system level, we average each metric over a system's summaries, overall and within each bucket, and ask whether the resulting order of the three systems matches the order of their human consistency rates. With three systems an order either matches or it does not; we therefore report the full table of scores rather than a rank correlation, which would take only a few distinct values. At the summary level, we ask how well a metric separates the summaries annotators judged consistent from those they judged inconsistent, independently of any threshold. Following the meta-evaluation practice of TRUE \citep{honovich2022true}, we report the area under the ROC curve of each metric against the adjudicated human label, pooled over the three systems, overall and within each bucket. An AUC of one half means the metric orders a consistent and an inconsistent summary at random; one means it always ranks the consistent summary higher.

Two quantities that look comparable are not. A metric score is not a probability that a summary is consistent, so we never compare a SummaC or QAFactEval score with a human consistency rate as if they were on one scale; we compare orders and discrimination only. Every number in the tables and in the text is computed by a script from the result files of the evaluation.
